\documentclass[reqno]{amsart}

\usepackage{amssymb}
\usepackage{microtype}

\newtheorem{thm}{Theorem}[section]
\newtheorem{prop}[thm]{Proposition}
\theoremstyle{definition}
\newtheorem{defn}[thm]{Definition}

\numberwithin{equation}{section}

\usepackage[hidelinks]{hyperref}

\begin{document}

\title{An introduction to no-arbitrage pricing of options contracts}
\author{John Myers}

\begin{abstract}
Later
\end{abstract}

\maketitle

\tableofcontents

\section{The two-asset market and binomial pricing model}

We suppose that the market contains two assets, one risky and the other risk-free, that we refer to as a \textit{stock} and \textit{bond}, respectively. We assume that the price process $S_t$ of the stock is stochastic, reflecting its risk, while we suppose that the bond carries a per-period interest rate $r>0$. This means that $\$1$ invested in the bond will increase to a value of $\$R$ one time period later, where $R = 1+r$. The rate $r$ is called the \textit{risk-free rate}, while $R$ is called the \textit{risk-free return}.

We imagine an investor owns a two-asset portfolio consisting of positions in the stock and the bond. The portfolio is held over a discrete time interval $t=0,1,\ldots,T$, beginning at $t=0$ and extending out to some time horizon $T>0$. At a given time $t$, we let $B_t$ be the value of the position in the bond, and we let $\Delta_t$ be the number of units of stock held in the portfolio. The total value of the portfolio at time $t$ is then given by the process
	\[
	V_t \stackrel{\text{def}}{=} B_t + S_t \Delta_t.
	\]
We allow $B_t$ and $\Delta_t$ to take negative values, corresponding to short positions. 

As time ticks forward, we assume that the investor adjusts their positions in the stock and bond, so $B_t$ and $\Delta_t$ are not constant processes. In fact, at a fixed time $t$, we assume that the investor will adjust \textit{both} their positions based on the price trajectory of the stock up to time $t$. This means that $B_t$ and $\Delta_t$ are themselves stochastic processes, special types that are \textit{adapted} to the natural filtration of the price process $S_t$. In other words, for each time $t$, the random variables $B_t$ and $\Delta_t$ are measurable with respect to the $\sigma$-algebra
	\[
	\mathcal{F}_t \stackrel{\text{def}}{=} \sigma(S_0,S_1,\ldots,S_t)
	\]
generated by the random variables $S_0,S_1,\ldots,S_t$.

The discussion so far has been completely general, with no further assumptions on the price process $S_t$ other than its stochasticity. Now, however, we introduce a simple model for $S_t$:

\begin{defn}\label{defn:binomial}
Let $u$, $d$, $p$ and $x_0$ be real numbers, with
	\[
	0 < d < 1 < u \quad \text{and} \quad 0 < p < 1,
	\]
and let $T$ be a positive integer. A \textit{binomial model}, with parameters $(u,d,p,x_0,T)$, is a probability space $(\Omega, \mathcal{F},P)$ of the following form:
\begin{enumerate}
\item The sample space $\Omega$ consists of all $(T+1)$-tuples of the form
	\[
	\omega = (x_0, x_1,\ldots,x_T)
	\]
such that each quotient $x_{t+1} / x_t$ is either $u$ or $d$ (for $t=0,1,\ldots, T-1$). The number of times that $u$ appears as a quotient is called the number of \textit{up moves} in the sample point $\omega$, while the number of times that $d$ appears is called the number of \textit{down moves}.
\item The $\sigma$-algebra $\mathcal{F}$ is the power set on $\Omega$.
\item The probability measure $P$ is given by
	\begin{equation}\label{eq:real-world-measure}
	P \left( \{\omega\} \right) = p^n (1-p)^{T-n},
	\end{equation}
	where $n$ is the number of up moves in $\omega$.
\end{enumerate}
The numbers $u$ and $d$ are called the \textit{up-factor} and \textit{down-factor}, $p$ and $P$ are called the \textit{real-world probability} and \textit{real-world measure}, and $T$ is called the \textit{time horizon} of the model.
\end{defn}

Given a binomial model of the above form, there are $T+1$ coordinate projection random variables
	\begin{equation}\label{eq:binomial-proj}
	X_t: \Omega \to \mathbb{R}, \quad X_t\left( (x_0, x_1,\ldots,x_T)\right) = x_t,
	\end{equation}
with $X_0$ a constant random variable. Note that the range of $X_t$ is
	\[
	\{ u^n d^{t-n} x_0 : 0 \leq n \leq t\},
	\]
and that, by a simple counting argument, we have
	\begin{equation}\label{eq:binom-prob}
	P\left(X_t = u^nd^{t-n} x_0 \right) = \binom{t}{n} p^n (1-p)^{t-n}.
	\end{equation}
The presence of the binomial coefficients explains the name of the model.

Though technically a binomial model is a probability space, according to our definition in this note, very often we shall abuse the terminology and refer to the stochastic process $X_t$ as a binomial model.

For each $t$ with $0 \leq t \leq T$, set
	\[
	\mathcal{F}_t = \sigma(X_0,X_1,\ldots,X_t).
	\]
These $\sigma$-algebras define the \textit{natural filtration} of the binomial model. Given a sample point $\omega\in \Omega$, its length-$t$ \textit{partial trajectory} is the sequence
	\[
	(X_0(\omega), X_1(\omega),\ldots,X_t(\omega)).
	\]
The next result identifies the atoms of $\mathcal{F}_t$ and computes their probabilities in terms of these partial trajectories. We leave the simple proofs to the reader.

\begin{prop}\label{prop:atom-prob}
Let $(\Omega, \mathcal{F},P)$ be a binomial model with time horizon $T$, and let $A$ be an atom of $\mathcal{F}_t$.
\begin{enumerate}
\item Two sample points $\omega$ and $\omega'$ are both in $A$ if and only if their length-$t$ partial trajectories are equal.
\item Let $m$ be the number of up moves in a partial trajectory of a sample point $\omega \in A$, which is independent of the choice of $\omega$ by \textnormal{(1)}. We have $P(A) = p^m (1-p)^{t-m}$, where $p$ is the real-world probability of the model.
\end{enumerate}
\end{prop}

One of the most important and useful properties of the binomial model is given in:

\begin{thm}
Every binomial model is a Markov process.
\begin{proof}
Let $X_t$ be the stochastic process of a binomial model with time horizon $T$. It will suffice to show that
	\[
	P\left(X_{t+1} = x_{t+1} \mid X_1=x_1, \ldots,X_t=x_t\right) = P\left(X_{t+1} = x_{t+1} \mid X_t=x_t\right)
	\]
for all $t$ with $0 \leq t < T$ and all vectors $(x_1,\ldots,x_t)$ of the form
	\[
	(x_1,\ldots,x_t) = (X_1(\omega),\ldots,X_t(\omega))
	\]
for some point $\omega$ in the sample space $\Omega$ of the model. To do this, note that
	\[
	P\left(X_{t+1} = x_{t+1} \mid X_1=x_1, \ldots,X_t=x_t\right) = \frac{P(X_1=x_1, \ldots,X_t=x_t,X_{t+1}=x_{t+1})}{P(X_1=x_1, \ldots,X_t=x_t)}
	\]
will be $0$ except in the cases that $x_{t+1}/x_t$ is either $u$ or $d$. In the first case, the probability in the numerator is $p$ times the probability in the denominator, while in the second case it is $1-p$ times the probability in the denominator. (Both of these computations use Proposition \ref{prop:atom-prob}.) Thus the conditional probability is either $0$, $p$, or $1-p$.

On the other hand, we have
	\[
	P\left(X_{t+1} = x_{t+1} \mid X_t=x_t\right) = \frac{P(X_t=x_t, X_{t+1}=x_{t+1})}{P(X_t=x_t)}.
	\]
Again, this probability will be $0$ except in the cases that $x_{t+1}/x_t$ is either $u$ or $d$. In both of these latter cases, the probability of the denominator is
	\[
	P(X_t=x_t) = \binom{t}{n} p^n (1-p)^{t-n},
	\]
for some $n$ with $0\leq n \leq t$, from \eqref{eq:binom-prob}. But the same counting argument that establishes \eqref{eq:binom-prob} also shows that the probability of the numerator is either
	\[
	\binom{t}{n} p^{n+1} (1-p)^{t-n} \quad \text{or} \quad \binom{t}{n} p^{n} (1-p)^{t-n+1},
	\]
depending on whether $x_{t+1}/x_t$ is $u$ or $d$. The desired result then follows immediately.
\end{proof}
\end{thm}

We now put this machinery to use. We assume that the price process $S_t$ of the stock, introduced at the beginning of this section, is the coordinate process $X_t$ of a binomial model with parameters $(u,d,p,S_0,T)$, in the sense of Definition \ref{defn:binomial}.

Returning to our two-asset market, we should expect that the potential return for the stock exceeds the guaranteed return for the bond since most investors are (at least mildly) risk-averse. Assuming that the stock price follows a binomial model, this means that we should have $R < u$, where $R$ is the risk-free return of the bond and $u$ is the up-factor of the binomial model. However, this inequality also follows from requiring the market to be free of \textit{arbitrage}. In finance, this latter word generally refers to the potential for a risk-free profit, and mathematical market and pricing models are designed in order to avoid them.

An arbitrage appears in the two-asset market if the price process of the stock follows a binomial model with $R \geq u$. Indeed, if this were the case, the investor could short sell a unit of the stock at time $t=0$ and invest the cash into the bond; hence $\Delta_0=-1$ and $B_0 = S_0$. Then, the initial value of the portfolio is
	\[
	V_0 = B_0 + S_0\Delta_0 = S_0 - S_0 \cdot 1 =0,
	\]
while one time period later the value of the bond position increases to
	\[
	RB_0 = RS_0 \geq uS_0 \geq S_1.
	\]
Thus, at time $t=1$, the investor has enough cash in the bond to cover the short position in the stock and pocket the remaining $RB_0 - S_1 \geq 0$ as profit. Since the portfolio does not require any (net) initial outlay, this is an arbitrage.

So, in order to avoid arbitrage, we shall assume $R < u$. But $d<1$ and $R>1$, so we have the pair of inequalities
	\[
	d < R < u.
	\]
We shall refer to these inequalities as the \textit{no-arbitrage condition}.

Finally, note that the stock and bond were completely independent assets until we specified that the price of the stock followed a binomial model and introduced the no-arbitrage condition: This condition mixes $u$ and $d$, which are parameters of $S_t$, with $R$, which is a parameter of the bond. Therefore, because it expresses a link between the assets, the no-arbitrage condition is really a property of the market, and accordingly the market is said to be \textit{arbitrage-free} if the no-arbitrage conditions hold.

\section{Arbitrage and martingales}

The no-arbitrage conditions guarantee that the number
	\[
	q \stackrel{\text{def}}{=} \frac{R-d}{u-d}
	\]
satisfies $0 < q < 1$. Using $q$, we may define a second probability measure $Q$ on the underlying measurable space $(\Omega, \mathcal{F})$ of the binomial model, by analogy with \eqref{eq:real-world-measure}, via
	\[
	Q(\{\omega\}) \stackrel{\text{def}}{=} q^n (1-q)^{T-n}
	\]
for each $\omega \in \Omega$, where $n$ is the number of up moves in $\omega$. Notice that this is the same formula \eqref{eq:real-world-measure} for the real-world measure $P$, but with $q$ in place of $p$.

This measure $Q$ is called the \textit{risk-neutral measure} of the model, terminology that will be explained in Section \ref{sec:claim-pricing}. Its name contrasts with the name of $P$, which is supposed to model the \textit{actual} probability of up and down moves of the stock price. Relative to $Q$, the probability of an up move is $q$, while the probability of a down move is
	\[
	1-q = \frac{u-R}{u-d}.
	\]
Notice that these two probabilities do not involve $P$ whatsoever, and so they have no relation to the real-world price dynamics of the stock. Rather, they are simply the (normalized) distances between the risk-free return $R$ and the up- and down-factors of the model: The risk-neutral probability of an up move is greater than that of a down move, or $q> 1-q$, when the risk-free return is closer to the up-factor than to the down-factor, and similarly for the reverse.

Though the probabilities they assign mean quite different things, the real-world and risk-neutral measures do have one feature in common, however: They agree on which events in $\Omega$ have probability $0$, which means the two measures are \textit{equivalent}.\footnote{In measure-theoretic language, two measures are \textit{equivalent} if each one is absolutely continuous with respect to the other.} Indeed, since $0<p<1$ and $0<q<1$, every singleton $\{\omega\}$ has strictly positive probability under both $P$ and $Q$, and hence so does every nonempty $A \subset \Omega$, being a union of such singletons. Thus $P(A)=0$ if and only if $A=\emptyset$, and similarly for $Q$, so the two measures are equivalent.

So, the absence of arbitrage in the two-asset market---with the stock price $S_t$ following a binomial model---implies the existence of the risk-neutral measure $Q$. In turn, as we will see below, the risk-neutral measure has the special feature that the \textit{discounted price process} $\tilde{S}_t$, defined by 
	\[
	\tilde{S}_t = \frac{S_t}{R^t}
	\]
for each $t=0,1,\ldots,T$, is a martingale when expectations are computed relative to $Q$. For this reason, the risk-neutral measure is also called an \textit{equivalent martingale measure} (or \textit{EMM}), the \textit{equivalent} part in the name referring to the equivalence with $P$ mentioned above. This means that absence of arbitrage in the market implies the existence of an EMM, which is the first half of the First Fundamental Theorem of Asset Pricing. The other half is the converse of this statement.

We will prove the First Fundamental Theorem in the current context of the binomial pricing model below. But the proof requires a few basic facts about $\sigma$-algebras, random variables and conditional expectations on finite sample spaces, which we recall first.

Let $\mathcal{F}$ be a $\sigma$-algebra on a finite sample space $\Omega$.
\begin{enumerate}
\item The $\sigma$-algebra $\mathcal{F}$ has a particularly simple description in terms of its \textit{atomic decomposition}. Indeed, an \textit{atom} of  $\mathcal{F}$ is a minimal non-empty set $A$ in $\mathcal{F}$. This means that if $B$ is in $\mathcal{F}$ and $B\subset A$, then either $B=\emptyset$ or $B=A$. One may show that the atoms of $\mathcal{F}$ partition the sample space $\Omega$, and that all non-empty sets in $\mathcal{F}$ are unions of atoms. Thus, the $\sigma$-algebra $\mathcal{F}$ determines---and is determined by---its atoms.

\item A function $X: \Omega \to \mathbb{R}$ is $\mathcal{F}$-measurable (i.e., it is a random variable) if and only if it is constant on the atoms of $\mathcal{F}$.

\item Atomic decompositions give a convenient and simple formula for the expectation $E(X \mid \mathcal{G})$ of an $\mathcal{F}$-measurable random variable $X: \Omega \to \mathbb{R}$, conditioned on a sub-$\sigma$-algebra $\mathcal{G}$ of $\mathcal{F}$. To make sense of this formula, we recall that the conditional expectation is a particular $\mathcal{G}$-measurable random variable, and so it must be constant on the atoms of $\mathcal{G}$. If $B$ is such an atom, then we have $B\in \mathcal{F}$ as well, and so it may be written (uniquely) as a union $B = \bigcup_{i\in I} A_i$, where the $A_i$ are atoms of $\mathcal{F}$. Then the desired formula is
	\[
	E(X \mid \mathcal{G})(\omega) = \frac{\sum_{i\in I} X(A_i)P(A_i)}{P(B)}
	\]
for every $\omega \in B$, where we write $X(A_i)$ for the unique value of $X$ on the atom $A_i$ (this uses $\mathcal{F}$-measurability of $X$). In the special case that $\mathcal{F}$ is the power set on $\Omega$ (as it will be, below), this formula may be written as
	\[
	E(X \mid \mathcal{G})(\omega) = \frac{\sum_{\eta\in B} X(\eta)P(\{\eta\})}{P(B)},
	\]
for every $\omega \in B$.
\end{enumerate}

If $X_t$ is a stochastic process, indexed $t=0,1,\ldots,T$, we recall from our earlier discussion that its \textit{natural filtration} consists of the $\sigma$-algebras
	\[
	\mathcal{F}_t \stackrel{\text{def}}{=} \sigma(X_0,X_1,\ldots,X_t),
	\]
one for each $t$. The atoms of $\mathcal{F}_t$ are exactly the inverse images of the random vector
	\[
	\omega \mapsto (X_0(\omega), X_1(\omega), \ldots, X_t(\omega)).
	\]
Using the ``atomic criterion'' for measurability mentioned above, the proof of the next proposition is trivial and is left to the reader.
	
\begin{prop}
Let $(\Omega, \mathcal{F},P)$ be a binomial model, as described in Definition \ref{defn:binomial}, and let $\mathcal{F}_t$ be the natural filtration of the stochastic process $Y_t$ consisting of the natural projection random variables given in \eqref{eq:binomial-proj}. Then a function $X: \Omega \to \mathbb{R}$ is $\mathcal{F}_t$-measurable if and only if it has the following property: Whenever $\omega$ and $\omega'$ are sample points in $\Omega$ such that
	\[
	(Y_0(\omega), Y_1(\omega), \ldots, Y_t(\omega)) = (Y_0(\omega'), Y_1(\omega'), \ldots, Y_t(\omega')),
	\] 
then $X(\omega) = X(\omega')$. \qedhere
\end{prop}	
	
Remember that we require the processes $B_t$ and $\Delta_t$ to be $S_t$-adapted, which just means that both $B_t$ and $\Delta_t$ are $\mathcal{F}_t$-measurable, for all $t$.

\begin{thm}[First Fundamental Theorem of Asset Pricing --- Binomial Pricing Model]
Consider a two-asset market consisting of a risky stock and a risk-less bond with risk-free return $R$, as described above, and suppose that the price process $S_t$ of the stock follows a binomial pricing model with time horizon $T>0$ and natural filtration $\mathcal{F}_t$. Then the market is arbitrage-free if and only if there exists a measure $Q$, equivalent to $P$, for which the discounted price process of the stock is a martingale:
	\[
	\tilde{S}_t = E_Q ( \tilde{S}_r\mid \mathcal{F}_t),
	\]
for all $t$ and $r$ with $0 \leq t \leq r\leq T$.
\begin{proof}
We first suppose that the market is arbitrage free, and we let $Q$ be the risk-neutral measure. We claim that the discounted price process is a martingale with respect to $Q$, and to prove this it will suffice to show
	\[
	RS_t(\omega) = E_Q\left(S_{t+1} \mid \mathcal{F}_t \right)(\omega)
	\]
for each $t=0,1,\ldots,T-1$ and for all $\omega\in \Omega$, where $\Omega$ is the sample space of the binomial model.

Fix $\omega$ and let $A$ be the atom of $\mathcal{F}_t$ containing it. Using the risk-neutral probability $q$, we may write
	\[
	Q(\{\omega\}) = \pi q^n (1-q)^{T-t-n}
	\]
for some integer $n$ and real number $\pi$ with $0\leq n \leq T-t$ and $0 < \pi \leq 1$. Let $A^u$ and $A^d$ consist of those $\eta \in A$ for which
	\[S_{t+1}(\eta) = uS_t(\omega) \quad \text{and} \quad S_{t+1}(\eta) = d S_t(\omega),
	\]
respectively. Then we have
	\[
	Q(A) = \pi, \quad Q(A^u) = \pi q, \quad \text{and} \quad Q(A^d) = \pi (1-q).
	\]
We now compute:
	\begin{align*}
	E\left(S_{t+1} \mid \mathcal{F}_t \right)(\omega) &= \frac{\sum_{\eta\in A} S_{t+1}(\eta)Q(\{\eta\})}{Q(A)} \\
	&= \frac{\sum_{\eta\in A^u} S_{t+1}(\eta)Q(\{\eta\}) + \sum_{\eta\in A^d} S_{t+1}(\eta)Q(\{\eta\})}{\pi} \\
	&= \left[ \frac{\sum_{\eta\in A^u} uQ(\{\eta\}) + \sum_{\eta\in A^d} dQ(\{\eta\})}{\pi} \right]S_t(\omega) \\
	&= \left[ \frac{\pi puq + \pi pd(1-q)}{\pi} \right]S_t(\omega) \\
	&= \left[ qu + (1-q)d\right]S_t(\omega) \\
	&= R S_t(\omega),
	\end{align*}
as desired.

For the converse, we suppose that $Q$ is an equivalent measure for which the discounted price process is a martingale. Let $\Omega^u$ and $\Omega^d$ be the two atoms of $\mathcal{F}_1 = \sigma(S_1)$, the first corresponding to those $\omega\in \Omega$ for which $S_1(\omega) = uS_0$, and the latter those $\omega$ for which $S_1(\omega) = dS_0$. Then, we compute:
	\begin{align*}
	RS_0 &= E_Q\left(S_1\right) \\
	&= \sum_{\omega\in \Omega} S_1(\omega) Q(\{\omega\}) \\
	&=\sum_{\omega\in \Omega^u} S_1(\omega) Q(\{\omega\}) + \sum_{\omega\in \Omega^d} S_1(\omega) Q(\{\omega\}) \\
	&= uS_0Q(\Omega^u) + dS_0Q(\Omega^d),
	\end{align*}
and so
	\[
	R = uQ(\Omega^u) + dQ(\Omega^d).
	\]
But $Q(\Omega^d) = 1 - Q(\Omega^u)$, from which it follows that $d < R < u$. 
\end{proof}
\end{thm}

In fact, an EMM has more ``martingale properties'' than just the one identified in the First Fundamental Theorem. One more such property involves the value process $V_t$ of the two-asset portfolio held by our investor, defined by
	\[
	V_t = B_t + S_t \Delta_t,
	\]
for each $t=0,1,\ldots,T$. Just as the price process $S_t$ has its discounted cousin, so too does the value process have a discounted version defined
	\[
	\tilde{V}_t \stackrel{\text{def}}{=} \frac{V_t}{R^t},
	\]
for each $t=0,1,\ldots,T$. The second martingale property possessed by $Q$ states that $\tilde{V}_t$ is a martingale when all expectations are computed relative to $Q$, provided that the two-asset portfolio is \textit{self-financing}. This latter terminology just means that, for each $t=0,1,\ldots,T-1$, we have an equation
	\[
	V_{t+1} = RB_t + S_{t+1}\Delta_t.
	\]
We understand these equations as follows: Exactly at time $t$, the investor adjusts their bond and stock positions, resulting with $\$B_t$ invested in the bond and $\Delta_t$ many units of stock. These positions are held over the interval $[t,t+1]$, resulting in a portfolio valued at
	\begin{equation}\label{eq:old-portfolio}
	RB_t + S_{t+1}\Delta_t
	\end{equation}
at time $t+1$. The investor then adjusts their positions again, ending with $\$B_{t+1}$ invested in the bond and $\Delta_{t+1}$ many units of stock, resulting in a portfolio valued at
	\begin{equation}\label{eq:new-portfolio}
	B_{t+1} + S_{t+1} \Delta_{t+1}
	\end{equation}
at time $t+1$. The self-financing condition requires that, however the investor adjusts their positions at time $t+1$, the total value \eqref{eq:new-portfolio} of the ``new'' portfolio is equal to the value \eqref{eq:old-portfolio} of the ``old'' portfolio held over the interval $[t,t+1]$. No capital infusions or withdrawals are required to build the new portfolio---in other words, it is self-financing.

Observe that the proof of the next martingale property of an EMM \textit{only} uses that the the discounted price process is a martingale, it does not use anything specific to a binomial model.

\begin{thm}\label{thm:value-process-martingale}
Consider a two-asset market consisting of a risky stock and a risk-less bond with risk-free return $R$, as described above, let $S_t$ be the price process of the stock with time horizon $T>0$, and let $Q$ be an EMM. Provided that the value process $V_t$ of a portfolio is self-financing and expectations are computed relative to $Q$, the discounted value process is a martingale with respect to the natural filtration $\mathcal{F}_t$ of $S_t$:
	\[
	\tilde{V}_t = E_Q ( \tilde{V}_r \mid \mathcal{F}_t )
	\]
for $0 \leq t \leq r \leq T$.

\begin{proof}
It will suffice to prove
	\[
	RV_t = E_Q\left( V_{t+1} \mid \mathcal{F}_t \right)
	\]
for each $t=0,1,\ldots,T-1$. Since the portfolio is self-financing, we have
	\[
	V_{t+1} = R B_t + S_{t+1} \Delta_t.
	\]
We then compute:
	\begin{align*}
	E_Q\left( V_{t+1} \mid \mathcal{F}_t \right) &= RE_Q(B_t \mid \mathcal{F}_t) + E_Q(S_{t+1}\Delta_t \mid \mathcal{F}_t) \\
	&= RB_t + E_Q(S_{t+1} \mid \mathcal{F}_t) \Delta_t \\
	&= RB_t + RS_t \Delta_t \\
	&= R V_t,
	\end{align*}
where we used $\mathcal{F}_t$-measurability of $B_t$ and $\Delta_t$ in going from the first line to the second, and that $Q$ is an equivalent martingale measure in going from the second to the third.
\end{proof}
\end{thm}

\section{Contingent claims and no-arbitrage pricing}\label{sec:claim-pricing}

The most general types of option contracts (here called \textit{contingent claims}) are defined as follows:

\begin{defn}
Let $S_t$ be the price process of the stock, indexed $t=0,1,\ldots,T$. Then a \textit{contingent claim}, with \textit{maturity} $T$, is an $S$-adapted process $\Phi_t$.

If $\Phi_t=0$ for all $t$ with $0 \leq t < T$, then the claim is called a \textit{non-early-exercise} claim, and the claim is said to \textit{not allow early exercise}. Otherwise, it is called an \textit{early-exercise} claim, and the claim is said to \textit{allow early exercise}.

The stock is called the \textit{underlying asset} of the claim, or simply the \textit{underlying}.
\end{defn}

The primary types of contingent claims discussed in this note are \textit{call} and \textit{put options}, of three different flavors.

\begin{defn}\label{defn:options}
Let $S_t$ be the price process of the stock, indexed $t=0,1,\ldots,T$, and let $K$ be a positive real number, called a \textit{strike price}.
\begin{enumerate}
\item The \textit{European call option}, with strike $K$ and maturity $T$, is given by
	\[
	\Phi^\text{Euro-call}_t = \begin{cases} 0 & : 0 \leq t < T, \\ S_T-K &: t=T.\end{cases}
	\]
The \textit{European put option}, with strike $K$ and maturity $T$, is given by
	\[
	\Phi^\text{Euro-put}_t = \begin{cases} 0 & : 0 \leq t < T, \\ K-S_T &: t=T.\end{cases}
	\]
\item The \textit{Asian call option}, with strike $K$ and maturity $T$, is given by
	\[
	\Phi^\text{Asian-call}_t = \begin{cases} 0 & : 0 \leq t < T, \\ \overline{S}-K &: t=T.\end{cases}
	\]
The \textit{Asian put option}, with strike $K$ and maturity $T$, is given by
	\[
	\Phi^\text{Asian-put}_t = \begin{cases} 0 & : 0 \leq t < T, \\  K-\overline{S} &: t=T.\end{cases}
	\]
In both types, we write $\overline{S}$ for the average stock price, given by
	\[
	\overline{S} = \frac{S_0+S_1+\cdots + S_T}{T+1}.
	\]
\item The \textit{American call} and \textit{put options}, both with strike $K$ and maturity $T$, are given by
	\[
	\Phi^\text{American-call}_t = S_t-K \quad \text{and} \quad \Phi^\text{American-put}_t = K-S_t
	\]
for $0\leq t \leq T$.
\end{enumerate}
\end{defn}

European and Asian options do not allow early exercise, while the American ones do.

To understand the real-world, financial intuition behind these types of claims, we will use a European call option on the stock as an example: At some initial time $t=0$, two parties write and agree on a contract that gives one of them (the \textit{option holder}, or \textit{buyer}) the right, but not the obligation, to purchase a number of units of stock from the other party (the \textit{option writer}, or \textit{seller}) at some future time $t=T$ at some price $\$K$. Crucially, all of these parameters---the maturity and the strike, as well as the number of stocks to be purchased---are set in the contract at the initial time $t=0$.

Clearly, the only case in which the option holder would exercise their right to buy the stocks is if the actual market price $S_T$ of the stock at maturity $t=T$ (the \textit{spot price}) is greater than the strike price. If it is, then the option holder can immediately turn around and sell the stocks on the market at $\$S_T$ per share, pocketing $\$(S_T-K)$ per share in profit. If instead the strike exceeds the spot price, the option holder will choose \textit{not} to exercise their right to buy, and the contract will expire without any exchange.

This describes a real-world European call option, and the abstract versions in Definition \ref{defn:options} function in a similar way. Indeed, we imagine that these latter contracts give the holder the right, but not the obligation, to \textit{exercise} at any time $t$ and collect $\$\Phi_t$ from the writer. We interpret $\Phi_t<0$ to mean that the \textit{holder owes the writer} an amount $\$(-\Phi_t)$, obviously an undesirable outcome from the holder's perspective, and so the holder would never choose to exercise at these times. Once the holder decides to exercise, the contract expires. The holder also has the option \textit{never} to exercise and to let the contract mature worthless---in this last case, we say that the holder ``exercises'' at time $t=\infty$.

Once the holder has a contingent claim, they need to solve an optimization problem: With the goal to maximize their profit, when  should they exercise? For the holder of a non-early-exercise claim, the optimal exercise times are obvious since the only possible time to turn a positive profit comes at maturity. Therefore, the holder should exercise at $t=T$, if $\Phi_T>0$. Otherwise, if $\Phi_T \leq 0$, the owner should let the contract mature worthless. Optimal exercise times for early-exercise claims take much more work to identify, which we save for Section \ref{sec:american}.

These contingent claims cost money. At the initial time $t=0$, the holder pays a premium to the writer to purchase the rights afforded to them in the contract. The basic problem is to identify a fair value for this premium.

Basic finance tells us that the price of an asset should be the sum of all present values of future cash flows, but if the asset is a contingent claim, we are immediately presented with three problems:
\begin{enumerate}
\item The holder can exercise at any time and collect their payoff. But this time will depend on the price trajectory of the stock, and hence \textit{the exercise time is random.}
\item Even if we knew the exercise time, \textit{the payoff itself is random}.
\item Even if we knew both the exercise time and the exact payoff, \textit{there is no obvious choice for the discount factor} with which to compute its present value.
\end{enumerate}
Of course, as we mentioned, the first is really only a problem for early-exercise claims. For European and Asian options and others like them, which do not allow early exercise, we can remove $\infty$ from the set of allowed exercise times and force exercise at maturity $t=T$, while still keeping the same cash flows, by replacing $\Phi_T$ with
	\begin{equation}\label{eq:zero-trick}
	\Phi_T^+ \stackrel{\text{def}}{=} \max{\left(\Phi_T, 0\right)}.
	\end{equation}
For a European call, this means that we consider
	\[
	\Phi_T^+ = \max{\left(S_T - K, 0 \right)}.
	\]
If $S_T< K$, then exercise at maturity $t=T$ results in the holder collecting $\$0$ from their counterparty, which is the same as the original European call with exercise time $t=\infty$.

\textit{In this first section on contingent claims, we shall only consider pricing non-early-exercise claims}, and leave the more complicated ones to a later section.

Now, one approach that attempts to overcome the second and third problems identified above is to embed the underlying asset into a market along with a risk-less asset---a market identical to the two-asset one we have been discussing so far---and to appeal to a no-arbitrage argument. Indeed, consider a (non-early-exercise) contingent claim $\Phi_T$ on our risky stock. Whatever its fair price $X_0$ at time $t=0$ turns out to be, imagine that it is possible to construct a portfolio, including positions in the bond and the stock, in such a way that the final value
	\[
	RB_{T-1} + S_{T} \Delta_{T-1}
	\]
of the portfolio held from $[T-1, T]$ coincides exactly with the final payoff $\Phi_T^+$. Then, provided that the portfolio is also self-financing, an investor who purchases such a portfolio makes a \textit{single} outlay of capital at time $t=0$ to purchase the portfolio, and ends at time $t=T$ with a portfolio whose value perfectly replicates the claim. If the purchase price $X_0$ of the claim was not equal to the purchase price $V_0$ of the replicating portfolio, then an arbitrage occurs: The investor would only have to short sell the more expensive asset at time $t=0$ (either the claim or the replicating portfolio), buy the less expensive one with the cash from the sale, pocketing the difference as profit, and then cover the short position at time $t=T$ with the payoff.

We are thus led to the following no-arbitrage pricing scheme: Given a (non-early-exercise) contingent claim $\Phi_T$, construct a self-financing portfolio that perfectly replicates the payoff $\Phi_T^+$ of the claim. Then, provided that such a portfolio is unique, the purchase price $X_0$ of the claim is taken to be the initial value $V_0$ of the replicating portfolio.

In order for this pricing strategy to work, we require both \textit{existence} and \textit{uniqueness} of these replicating portfolios. In the case that the price process for the underlying stock is a binomial model, these requirements are easily handled:

\begin{thm}
Consider a two-asset market consisting of a risky stock and a risk-less bond, as described above, and suppose that $S_t$ is the price process of the stock with time horizon $T>0$. Provided that $S_t$ is a binomial model, every non-early-exercise contingent claim $\Phi_T$ on the stock has a unique replicating portfolio.

\begin{proof}
The method of proof is a familiar ``backwards induction,'' as it is known in the literature. The idea is to assume that the desired processes have been constructed backwards
	\[
	B_{T-1},B_{T-2} \ldots, B_{t+1}, \quad \text{and} \quad \Delta_{T-1},\Delta_{T-2} \ldots, \Delta_{t+1},
	\]
beginning with $B_{T-1}$ and $\Delta_{T-1}$, and going back to some time $t$ with $0 \leq t \leq T-2$. Given this, the inductive step is then to show how to construct $B_t$ and $\Delta_t$.

Of course, before establishing the inductive step, one needs to handle the base case, which means constructing $B_{T-1}$ and $\Delta_{T-1}$. However, as we will show at the end of the proof, the argument for the base case is the same as the argument for the inductive step. So, we go straight to the proof of the inductive step.

Assuming $B_{t+1}$ and $\Delta_{t+1}$ exist, we have the portfolio value
	\[
	V_{t+1} = B_{t+1} + S_{t+1}\Delta_{t+1}.
	\]
The portfolio is supposed to be self-financing, and so we must have
	\begin{equation}\label{eq:self-financing}
	V_{t+1} = R B_t + S_{t+1}\Delta_t.
	\end{equation}
Note that this is an equation in the two unknowns $B_t$ and $\Delta_t$. We shall prove that it has a unique solution.

Since both $B_t$ and $\Delta_t$ are supposed to be $\mathcal{F}_t$-measurable, we need only define these random variables on the atoms of this $\sigma$-algebra. Let $A$ be such an atom, and let $\omega \in A$. The atom contains (at least) two trajectories, $\omega_u$ and $\omega_d$, with the properties that
	\[
S_{t+1}(\omega_u) = u S_{t}(\omega) \quad \text{and} \quad S_{t+1}(\omega_d) = d S_{t}(\omega),
	\]
where $u$ and $d$ are the up- and down-factors of the model. Then, we have
	\[
	V_{t+1}(\omega_u) = R B_{t}(\omega) + uS_{t}(\omega) \Delta_{t}(\omega)
	\]
and
	\[
	V_{t+1}(\omega_d) = R B_{t}(\omega) + dS_{t}(\omega) \Delta_{t}(\omega)
	\]
from \eqref{eq:self-financing}, in view of the equations
	\[B_t(\omega) = B_t(\omega_u) = B_t(\omega_d) \quad \text{and} \quad \Delta_t(\omega) = \Delta_t(\omega_u) = \Delta_t(\omega_d).
	\]
We then have the linear system
	\[
	\begin{bmatrix}
	R & uS_{t}(\omega) \\
	R & dS_{t}(\omega)
	\end{bmatrix}
	\begin{bmatrix} B_{t}(\omega) \\ \Delta_{t}(\omega) \end{bmatrix} =
	\begin{bmatrix}
	V_{t+1}(\omega_u) \\
	V_{t+1}(\omega_d)
	\end{bmatrix}
	\]
in the unknowns $B_{t}(\omega)$ and $\Delta_{t}(\omega)$. The determinant of the coefficient matrix is nonzero, and hence a unique solution exists. This completes the proof of the inductive step.

Now, we return to the base case, demonstrating how to construct $B_{T-1}$ and $\Delta_{T-1}$. The portfolio is supposed to replicate the contingent claim, and hence we must have
	\[
	\Phi_T^+ = RB_{T-1} + S_T \Delta_{T-1}.
	\]
This is an equation in the two unknowns $B_{T-1}$ and $\Delta_{T-1}$ which is of the same form as \eqref{eq:self-financing} if we set $V_T = S_T$ and take $t=T-1$. Thus, the argument establishing the base case is the same as the argument for the inductive step.
\end{proof}
\end{thm}

Provided that self-financing, replicating portfolios exist and are unique, the no-arbitrage argument was used to show that the initial price $X_0$ of a non-early-exercise contingent claim $\Phi_T$ on the stock coincides with the initial value $V_0$ of its replicating portfolio. It is natural to extend this to all times $t=0,1,\ldots,T$, defining the entire price process $X_t$ of the claim, by setting $X_t=V_t$, where we take $V_T = \Phi_T^+$. In the case that the price process of the stock is a binomial model, the no-arbitrage condition also provides the existence of a unique equivalent martingale measure $Q$, by the First Fundamental Theorem. From Theorem \ref{thm:value-process-martingale}, we then get concise formulas for the no-arbitrage price process of the claim:
	\begin{equation}\label{eq:martingale-pricing}
	X_t = \frac{1}{R^{r-t}} E_Q\left( X_r \mid \mathcal{F}_t\right),
	\end{equation}
for $0 \leq t \leq r \leq T$. In particular, the initial price is given by
	\begin{equation}\label{eq:expectation-pricing}
	X_0 = \frac{1}{R^T} E_Q\left( \Phi_T^+ \right).
	\end{equation}
This latter formula is especially enlightening in view of the discussion at the beginning of this section, and how it navigates the second two problems listed there:
\begin{enumerate}
\item[(2)] The problem of not knowing the future payoff of the contingent claim with certainty, because it is a function of a random variable, is circumvented by considering the \textit{expected} future payoff. However, one must take care to note that \textit{the expected value is computed relative to the risk-neutral measure $Q$}, and not the real-world measure $P$.
\item[(3)] The problem of not knowing the correct discount factor is solved by using the risk-free rate. This is of interest since, in the real world, the correct discount factor would take into account the risk profile of the claim and the stock; however, in this formula, the risk-neutral measure $Q$ absorbs and reflects this risk, enabling one to discount at the risk-free rate.
\end{enumerate}
Since the formula \eqref{eq:expectation-pricing} prices the claim in terms of its expected payoff only, it is a \textit{risk-neutral} formula, which inspires the name of the risk-neutral measure $Q$.\footnote{An analogy is helpful. Suppose a player is selecting one of two games to play. In the first, the player wins or loses $\$100$ with equal probability $0.5$, while in the second, they win or lose $\$10$, also with equal probability. A \textit{risk-seeking} player would choose the first game and a shot at higher winnings, at the expense of risking higher losses, while a \textit{risk-averse} player would choose the latter, less risky game. But a \textit{risk-neutral} player would only consider the expected payoff of each game, which is $\$0$ for both, and would thus not express a preference for one game over the other. (Admittedly, the analogy with the risk-neutral measure is not perfect, because $Q$ \textit{does} take into account risk.)}

\section{Optimal stopping}



\section{Pricing a European option with a binomial model}

\section{Pricing an American option with a binomial model}\label{sec:american}

A European put option allows the holder to decide, at maturity time $T$, whether to exercise the option and collect the payoff $K-S_T$, or to let the contract expire worthless if the strike $K$ is less than the spot price $S_T$. The point is: The European contract stipulates that the holder can \textit{only} exercise at the maturity time $T$, not before. An \textit{American put option}, on the other hand, is identical to its European cousin in every respect, \textit{except} that it allows the holder to exercise at \textit{any} time $t$, up to and including the maturity time $T$. So the decisions facing the owner of an American put option are, at all times $t$ with $0 \leq t < T$, whether to exercise the option and collect
	\[
	\Phi(S_t) \stackrel{\text{def}}{=} \max(K-S_t,0),
	\]
or to keep the contract alive and continue to time $t+1$. The decisions are made in the hopes to maximize profit, of course.

For the sake of argument, let's suppose that the owner of the option has let the contract carry to time $t=T-1$ without exercise, one time period before maturation at time $t=T$. They must then decide whether to exercise immediately and collect $\Phi(S_{T-1})$, or to wait and collect $\Phi(S_T)$. In making this decision, the difficulty is two-fold:
\begin{enumerate}
\item \textit{While $\Phi(S_{T-1})$ is a known quantity at time $t=T-1$, the quantity $\Phi(S_T)$ is not.} The option owner cannot see into the future, even just one time period ahead.
\item Even if $\Phi(S_T)$ were a deterministic quantity, the payoffs occur at different times, one at the present time $t=T-1$ and the other at a future time $t=T$, and so \textit{proper discounting must be applied to obtain the present value of $\Phi(S_T)$.}
\end{enumerate}
We encourage the reader to return to the discussion at the beginning of Section \ref{sec:claim-pricing}, and notice how these two problems align exactly with the problems we discussed there. The lessons learned in that section suggest that---provided the price process of the stock is a binomial model and the market is arbitrage-free---we can solve the first of these problems by replacing the random quantity $\Phi(S_T)$ with its expectation $E_Q(\Phi(S_T))$ relative to the risk-neutral measure $Q$, and solve the second problem by using the discount factor $1/R$, where $R$ is the risk-free return. Thus, the owner should exercise at time $t=T-1$ if
	\begin{equation}\label{eq:american-intrinsic-continuation}
	\Phi(S_{T-1})(\omega) > \frac{1}{R} E_Q\left(\Phi(S_T)\mid \mathcal{F}_{T-1}\right)(\omega),
	\end{equation}
or defer exercise to time $t=T$ if this inequality is reversed. The quantity $\Phi(S_{T-1})$ on the left-hand side of this inequality is called the \textit{intrinsic value} of the option at time $t=T-1$, while the quantity on the right-hand side is called its \textit{continuation value}.

Suppose, for each $\omega\in \Omega$, that we define
	\[
	V_T = \max{ \left[ K-S_T, 0 \right]}
	\]
and then set
	\[
	V_t = \max{ \left[ K-S_t, \frac{1}{R} E_Q\left( V_{t+1} \mid \mathcal{F}_t \right)\right]}
	\]
for $0 \leq t < T$.

\appendix
\section{Probability theory}

\subsection{Probability spaces}

In the abstract mathematical theory, a \textit{sample space} is just a nonempty set $\Omega$. However, in real-world applications, a sample space is often conceptualized as the set of all outcomes of some ``random experiment.'' The elements of a sample space are called \textit{sample points}

\begin{defn}
Let $\Omega$ be a finite sample space. A \textit{$\sigma$-algebra} on $\Omega$ is a collection $\mathcal{F}$ of subsets of $\Omega$ with the following properties:
\begin{enumerate}
\item The sample space $\Omega$ itself is contained in $\mathcal{F}$.
\item If $U$ is a set in $\mathcal{F}$, then the (absolute) complement $U^c = \Omega \smallsetminus U$ is also in $\mathcal{F}$.
\item If both $U$ and $V$ are sets in $\mathcal{F}$, then the union $U\cup V$ is also in $\mathcal{F}$.
\end{enumerate}
A set in $\mathcal{F}$ is called an \textit{$\mathcal{F}$-measurable set}, or simply a \textit{measurable set} if the identity of the $\sigma$-algebra is clear. In the probabilistic context, measurable sets are called \textit{events}.
\end{defn}

\begin{defn}
A \textit{finite event space} is a pair $(\Omega, \mathcal{F})$ consisting of a finite sample space and $\mathcal{F}$ is a $\sigma$-algebra on $\Omega$.\footnote{Event spaces are called \textit{measurable spaces} in general measure theory.}
\end{defn}

\begin{defn}
Let $\mathcal{F}$ be a $\sigma$-algebra on a finite sample space.
\begin{enumerate}
\item A \textit{sub-$\sigma$-algebra} of $\mathcal{F}$ is a $\sigma$-algebra $\mathcal{G}$ on $\Omega$ with $\mathcal{G} \subset \mathcal{F}$.
\item A \textit{filtration} of $\mathcal{F}$ is a family $(\mathcal{F}_t)_{t\in T}$ of sub-$\sigma$-algebras of $\mathcal{F}$, indexed by a linearly-ordered set $T$, such that $s \leq t$ $\Rightarrow$ $\mathcal{F}_s \subset \mathcal{F}_t$ for all $s,t\in T$.
\item A \textit{filtered $\sigma$-algebra} on $\Omega$ is a pair $(\mathcal{F}, (\mathcal{F}_t)_{t\in T})$ consisting of a $\sigma$-algebra $\mathcal{F}$ on $\Omega$ and a filtration $(\mathcal{F}_t)_{t\in T}$ of $\mathcal{F}$.
\end{enumerate}
\end{defn}

\begin{defn}
An \textit{atom} of a $\sigma$-algebra $\mathcal{F}$ (on a finite sample space $\Omega$) is a set $A$ in $\mathcal{F}$ with the following two properties:
\begin{enumerate}
\item It is nonempty, $A\neq \emptyset$.
\item If $B$ is a nonempty set in $\mathcal{F}$ with $B\subset A$, then $A=B$.
\end{enumerate}
\end{defn}

\begin{thm}[Atomic Partition Theorem]
Let $\Omega$ be a finite sample space. We have the following two operations connecting $\sigma$-algebras on $\Omega$ to partitions:
\begin{enumerate}
\item Every $\sigma$-algebra $\mathcal{F}$ on $\Omega$ yields a partition of $\Omega$, where the sets in the partition are the atoms of $\mathcal{F}$.
\item Every partition of $\Omega$ yields a $\sigma$-algebra $\mathcal{G}$ on $\Omega$, where the sets in $\mathcal{G}$ are all unions of sets in the partition, along with $\emptyset$.
\end{enumerate}
These two operations are mutual inverses of each other, so that every $\sigma$-algebra is uniquely determined by a partition and every partition is uniquely determined by a $\sigma$-algebra.
\end{thm}

\begin{defn}
Let $\mathcal{F}$ be a $\sigma$-algebra on a finite sample space $\Omega$. A \textit{probability measure} on $\mathcal{F}$ is a function $P:\mathcal{F} \to [0,\infty)$ with the following properties:
\begin{enumerate}
\item We have $P(\Omega)=1$.
\item If $U$ and $V$ are two $\mathcal{F}$-measurable sets that are disjoint, then $P(U\cup V) = P(U) + P(V)$.
\end{enumerate}
\end{defn}

\begin{defn}
A \textit{finite probability space} is a triple $(\Omega, \mathcal{F}, P)$ consisting of a finite sample space $\Omega$, a $\sigma$-algebra $\mathcal{F}$ on $\Omega$, and a probability measure $P$ on $\mathcal{F}$.

A \textit{finite filtered probability space} is a quadruple $(\Omega, \mathcal{F}, (\mathcal{F}_t)_{t\in T}, P)$ where the triple $(\Omega, \mathcal{F}, P)$ is a finite probability space and $(\mathcal{F}_t)_{t\in T}$ is a filtration of $\mathcal{F}$.
\end{defn}

\subsection{Random vectors and variables}

\begin{defn}
Let $(\Omega, \mathcal{F})$ be a finite event space. A function $\mathbf{X}: \Omega \to \mathbb{R}^d$ is called \textit{$\mathcal{F}$-measurable} if the preimage
	\[
	\{\mathbf{X} = \mathbf{x}\} \stackrel{\text{def}}{=} \{ \omega \in \Omega : \mathbf{X}(\omega)=\mathbf{x}\}
	\]
is in $\mathcal{F}$, for each $\mathbf{x}\in \mathbb{R}^d$. If the identity of the $\sigma$-algebra is clear from context, \textit{$\mathcal{F}$-measurable} will be shortened to just \textit{measurable}.
\end{defn}

\begin{defn}
Let $(\Omega, \mathcal{F})$ be a finite event space.
\begin{enumerate}
\item A function $\mathbf{X}: \Omega \to \mathbb{R}^d$ is called a \textit{$d$-dimensional random vector} if it is $\mathcal{F}$-measurable.
\item A function $X: \Omega \to \mathbb{R}$ is a \textit{random variable} if it is a $1$-dimensional random vector.
\item A \textit{random vector} on a finite probability space $(\Omega, \mathcal{F}, P)$ is a random vector on the underlying finite event space $(\Omega, \mathcal{F})$.
\end{enumerate}
\end{defn}

\begin{thm}[Atomic Criterion for Measurability]\label{thm:atomic-criterion-measurable}
Let $(\Omega, \mathcal{F})$ be a finite event space. A function $\mathbf{X}: \Omega \to \mathbb{R}^d$ is $\mathcal{F}$-measurable (i.e., it is a random vector) if and only if it is constant on the atoms of $\mathcal{F}$.
\end{thm}

\begin{defn}
Let $\Omega$ be a finite sample space and $\mathbf{X}: \Omega \to \mathbb{R}^d$ a function. The \textit{$\sigma$-algebra generated by $\mathbf{X}$}, denoted $\sigma(\mathbf{X})$, has atoms equal to the preimages $\{\mathbf{X} = \mathbf{x}\}$, one for each vector $\mathbf{x}\in \mathbb{R}^d$ for which the preimage is nonempty.
\end{defn}

\begin{thm}[Functional Criterion for Measurability]\label{thm:functional-criterion-measurable}
Let $\Omega$ be a finite sample space and let $\mathbf{X}: \Omega \to \mathbb{R}^d$ be a function. Then a function $\mathbf{Y}: \Omega \to \mathbb{R}^e$ is $\sigma(\mathbf{X})$-measurable if and only if there is a function $h: \mathbb{R}^d \to \mathbb{R}^e$ such that $\mathbf{Y} = h(\mathbf{X})$.
\end{thm}

This is a key result, so it deserves elaboration. Supposing that $\mathbf{Y}$ is $\sigma(\mathbf{X})$-measurable, a function $h$ as in the theorem is constructed as follows. If $\mathbf{x}\in \mathbb{R}^d$ is a vector whose preimage $\{\mathbf{X} = \mathbf{x}\}$ is nonempty, we define
	\[
	h(\mathbf{x}) = \mathbf{Y}(\{\mathbf{X} = \mathbf{x}\}),
	\]
where the right-hand side is the unique value of $\mathbf{Y}$ on the preimage $\{\mathbf{X} = \mathbf{x}\}$. For those vectors $\mathbf{x}\in \mathbb{R}^d$ such that $\{\mathbf{X} = \mathbf{x}\}$ is empty, we define $h(\mathbf{x})$ to be an arbitrary but fixed value. Note that the only part of $h$ that is unique is its restriction to the range of $\mathbf{X}$.


\begin{thm}
Let $\Omega$ be a finite sample space and $\mathbf{X}: \Omega \to \mathbb{R}^d$ a function.
\begin{enumerate}
\item The function $\mathbf{X}$ is $\sigma(\mathbf{X})$-measurable.
\item If $\mathcal{F}$ is a $\sigma$-algebra on $\Omega$ for which $\mathbf{X}$ is $\mathcal{F}$-measurable, then $\sigma(\mathbf{X}) \subset \mathcal{F}$.
\end{enumerate}
\end{thm}

\subsection{Lebesgue integration and expectations}

Let $X$ be a random variable on a finite event space $(\Omega, \mathcal{F})$. Given an atom $A$ of $\mathcal{F}$, by Theorem \ref{thm:atomic-criterion-measurable} the value $X(\omega)$ is constant as $\omega$ ranges over $A$. We write $X(A)$ for this unique value. Recall that we write $\alpha(\mathcal{F})$ for the set of atoms of $\mathcal{F}$.

\begin{defn}
Let $X$ be a random variable on a finite probability space $(\Omega, \mathcal{F}, P)$. Given an event $U\in \mathcal{F}$, we define the \textit{Lebesgue integral} of $X$ over $U$, denoted $\int_U X \, dP$, to be the real number
	\[
	\int_U X \, dP = \sum_{A\in \alpha(\mathcal{F})} I_U(A)X(A)P(A),
	\]
where $I_U$ is the indicator function of $U$.
\end{defn}

\begin{defn}
Let $X$ be a random variable on a finite probability space $(\Omega, \mathcal{F}, P)$. The \textit{expected value} of $X$, denoted $E(X)$, is the real number
	\[
	E(X) = \int_\Omega X \, dP.
	\]
\end{defn}

\subsection{\texorpdfstring{$L^2$-spaces}{L2-spaces}, conditional expectations and probabilities}

Many constructs in measure-theoretic probability theory are clarified through the introduction of a geometric viewpoint. This geometry comes through the Hilbert space structure on sets of random variables that we introduce in:

\begin{defn}\label{defn:l2}
Letting $(\Omega, \mathcal{F}, P)$ be a finite probability space, we write $L^2(\Omega, \mathcal{F}, P)$ for the set of all $\mathcal{F}$-measurable random variables on $\Omega$.
\begin{enumerate}
\item We define the \textit{sum} of two random variables $X$ and $Y$ in $L^2(\Omega, \mathcal{F}, P)$, denoted $X+Y$, to be the $\mathcal{F}$-measurable random variable given by
	\[
	(X+Y)(\omega) = X(\omega) + Y(\omega),
	\]
for all $\omega \in \Omega$. We define the \textit{product} of $X$ and $Y$, denoted $XY$, to be the $\mathcal{F}$-measurable random variable given by
	\[
	(XY)(\omega) = X(\omega) Y(\omega),
	\] 
for all $\omega\in \Omega$. Given a scalar $c\in \mathbb{R}$, we define the \textit{scalar multiple} of $c$ and $X$, denoted $cX$, to be the $\mathcal{F}$-measurable random variable given by
	\[
	(cX)(\omega) = c X(\omega),
	\]
for all $\omega \in \Omega$. Finally, we define the \textit{zero random variable} to be the constant random variable that sends every $\omega\in \Omega$ to $0\in \mathbb{R}$.
\item We define the \textit{inner product} of two random variables $X$ and $Y$ in $L^2(\Omega, \mathcal{F}, P)$, denoted $\langle X, Y \rangle$, to be the real number given by
	\[
	\langle X, Y\rangle = \int_\Omega XY \, dP.
	\]
\end{enumerate}
\end{defn}

Our definition of $H = L^2(\Omega, \mathcal{F}, P)$ as a set of \textit{actual} random variables presents a difficulty of what it should mean for two random variables in $H$ to be \textit{equal}. For example, suppose that $N$ is a set in $\mathcal{F}$ with zero probability, $P(N)=0$, and suppose that $X\in H$ is a random variable which is equal to $0$ on the complement $N^c$, but takes a nonzero value on $N$. Then $X$ is not equal to the zero random variable in $H$, but yet
	\[
	\int_\Omega X \, dP=0.
	\]
As we will see in Theorem \ref{thm:hilbert} below, this will mean that the inner product given in Definition \ref{defn:l2} is \textit{not} positive definite.

However, notice that $X$ is \textit{almost surely equal} to the zero random variable, in the following sense:

\begin{defn}
Let $X$ and $Y$ be two random variables on a finite probability space $(\Omega, \mathcal{F}, P)$. Then $X$ and $Y$ are \textit{almost surely equal} if the set
	\[
	N = \{\omega \in \Omega \mid X(\omega) \neq Y(\omega) \}
	\]
has probability $0$.
\end{defn}

There are then two ways to leverage this notion of \textit{almost sure equality} in a proper definition of the space $H = L^2(\Omega, \mathcal{F},P)$. First, instead of defining $H$ to be a set of actual random variables, one can define it as the set of equivalence classes of random variables, where the equivalence relation is almost sure equality. Strictly speaking, this is the proper definition. Alternatively, one can keep $H$ as a set of random variables, and just understand that equality of random variables in $H$ really means almost sure equality. This latter road is the one we take.

It will be convenient to write $\alpha'(\mathcal{F})$ for the subset of $\alpha(\mathcal{F})$ consisting of all atoms of nonzero probability.

With this convention in force, we then have the main theorem:

\begin{thm}\label{thm:hilbert}
Let $(\Omega, \mathcal{F}, P)$ be a finite probability space and set $H = L^2(\Omega, \mathcal{F}, P)$.
\begin{enumerate}
\item Equipped with the sum, scalar multiplication, and zero random vector defined in Definition \ref{defn:l2}, the set $H$ is a finite-dimensional vector space.
\item Equipped with the vector space structure in (1) and the inner product defined in Definition \ref{defn:l2}, the set $H$ is a Hilbert space.
\item The Hilbert space $H$ has an orthogonal basis consisting of all indicator functions $I_A$, as $A$ ranges over the set $\alpha'(\mathcal{F})$ of all atoms of $\mathcal{F}$ with nonzero probability.
\end{enumerate}
\end{thm}

We compute an \textit{orthonormal} basis by normalizing the indicator functions in the orthogonal basis identified in the theorem. For this, we use the following computation:
	\[
	\| I_A \|^2 = \langle I_A, I_A \rangle = \int_\Omega I_A^2 \, dP = P(A).
	\]
Thus, the space $H$ has an orthonormal basis consisting of all scaled indicator functions
	\begin{equation}\label{eq:onb}
	\phi_A \stackrel{\text{def}}{=} \frac{I_A}{\sqrt{P(A)}},
	\end{equation}
one for each atom $A\in \alpha'(\mathcal{F})$.

Suppose that $K$ is a vector subspace of $H$. Given a random variable $X\in H$, the \textit{orthogonal projection of $X$ onto $K$} is the unique random variable $\Pi$ such that $\Pi\in K$ and $\langle \Pi-X, Y \rangle = 0$ for all $Y\in K$. If $K$ has an orthonormal basis $\{Y_i\}_{i\in I}$, then $\Pi$ is given by its \textit{generalized Fourier expansion}:
	\begin{equation}\label{eq:fourier}
	\Pi = \sum_{i\in I} \langle X, Y_i \rangle Y_i.
	\end{equation}
Important examples of subspaces are of the form $K = L^2(\Omega, \mathcal{G}, P|_{\mathcal{G}})$, where $\mathcal{G}$ is a sub-$\sigma$-algebra of $\mathcal{F}$ and $P|_{\mathcal{G}}: \mathcal{G} \to [0,\infty)$ is the restriction of $P$ to $\mathcal{G}$. Orthogonal projections onto subspaces of this type have special names:

\begin{defn}
Let $X$ be a random variable on a finite probability space $(\Omega, \mathcal{F}, P)$, and let $\mathcal{G}$ be a sub-$\sigma$-algebra of $\mathcal{F}$. The \textit{conditional expectation of $X$ given $\mathcal{G}$}, denoted $E(X \mid \mathcal{G})$, is the orthogonal projection of $X$ onto the subspace $L^2(\Omega, \mathcal{G}, P|_{\mathcal{G}})$.
\end{defn}

Using the orthonormal basis of $K=L^2(\Omega, \mathcal{G}, P|_{\mathcal{G}})$ consisting of scaled indicator functions \eqref{eq:onb}, conditional expectations may be computed using their Fourier expansions \eqref{eq:fourier}. To do so, let an atom $B\in \alpha'(\mathcal{G})$ be given and consider the basis random vector $\phi_B$. We have
	\[
	\langle X, \phi_B \rangle \phi_B = \frac{\int_B X \, dP}{P(B)}  I_B,
	\]
and so
	\[
	E(X \mid \mathcal{G}) = \sum_{B \in \alpha'(\mathcal{G})} \frac{\int_B X \, dP}{P(B)}  I_B.
	\]
	
Directly from the definition of a conditional expectation as an orthogonal projection, one proves:

\begin{thm}
Let $X$ be a random variable on a finite probability space $(\Omega, \mathcal{F}, P)$, and let $\mathcal{G}$ be a sub-$\sigma$-algebra of $\mathcal{F}$.
\begin{enumerate}
\item The conditional expectation $E(X\mid \mathcal{G})$ is $\mathcal{G}$-measurable and
	\[
	\int_U X \, dP = \int_U E(X\mid \mathcal{G}) \, dP
	\]
for all $U\in \mathcal{G}$.
\item If $Z$ is another random variable with the two properties of $E(X\mid \mathcal{G})$ listed in (1), then $Z$ and $E(X\mid \mathcal{G})$ are equal almost surely.
\end{enumerate}
\end{thm}

In some texts on measure-theoretic probability theory, the properties of conditional expectation listed in this latter theorem are taken as the \textit{definition} of conditional expectation.

The characterizing property of conditional expectation as an orthogonal projection makes some of its most useful features easy to establish. We collect the ones we will need most often:

\begin{thm}\label{thm:cond-exp-properties}
Let $(\Omega, \mathcal{F}, P)$ be a finite probability space, let $X, Y \in L^2(\Omega, \mathcal{F}, P)$, and let $\mathcal{G}$ be a sub-$\sigma$-algebra of $\mathcal{F}$.
\begin{enumerate}
\item (\textit{Linearity}) For all $a,b \in \mathbb{R}$,
	\[
	E(aX + bY \mid \mathcal{G}) = a\,E(X\mid \mathcal{G}) + b\,E(Y \mid \mathcal{G}).
	\]
\item (\textit{Taking out what is known}) If $Y$ is $\mathcal{G}$-measurable, then
	\[
	E(XY \mid \mathcal{G}) = Y E(X \mid \mathcal{G}).
	\]
\item (\textit{Tower property}) If $\mathcal{H}$ is a sub-$\sigma$-algebra of $\mathcal{G}$, then
	\[
	E\big(E(X\mid \mathcal{G}) \mid \mathcal{H}\big) = E(X \mid \mathcal{H}).
	\]
\item (\textit{Triviality}) If $\mathcal{G} = \{\emptyset, \Omega\}$ is the trivial $\sigma$-algebra, then $E(X \mid \mathcal{G})$ is the constant random variable equal to $E(X)$.
\item (\textit{Idempotence}) If $X$ is $\mathcal{G}$-measurable, then $E(X \mid \mathcal{G}) = X$.
\end{enumerate}
\end{thm}

Each of these properties reflects a familiar fact about orthogonal projections, transplanted to the setting of $H$. For instance, (5) is simply the statement that a vector already lying in a subspace is its own projection onto that subspace, while (3) reflects the fact that $L^2(\Omega,\mathcal{H},P|_{\mathcal{H}})$ is itself a subspace of $L^2(\Omega,\mathcal{G},P|_{\mathcal{G}})$ whenever $\mathcal{H} \subset \mathcal{G}$, so that projecting onto the smaller subspace can be done either directly, or in two stages by way of the larger one.

Suppose that $\mathcal{G}$ is of the special form $\mathcal{G} = \sigma(\mathbf{Y})$, where $\mathbf{Y}$ is a $d$-dimensional $\mathcal{F}$-measurable random vector. Then we write $E(X \mid \mathbf{Y})$ in place of $E(X\mid \sigma(\mathbf{Y}))$. This conditional expectation is $\sigma(\mathbf{Y})$-measurable and so by the Functional Criterion for Measurability (Theorem \ref{thm:functional-criterion-measurable}) there is a function
	\begin{equation}\label{eq:cond-exp-func}
	E(X \mid \mathbf{Y} = {?}) :\mathbb{R}^d \to \mathbb{R}, \quad \mathbf{y} \mapsto E(X \mid \mathbf{Y} = \mathbf{y}),
	\end{equation}
such that
	\[
	E(X \mid \mathbf{Y} = \mathbf{y}) = E(X \mid \mathbf{Y})(\{ \mathbf{Y} = \mathbf{y} \})
	\]
for each $\mathbf{y} \in \mathbb{R}^d$ for which $P(\mathbf{Y} = \mathbf{y})>0$. If this probability is $0$, then $E(X\mid \mathbf{Y}=\mathbf{y})=0$. We will show later in Theorem \ref{thm:cond-exp-familiar} that the function $E(X\mid \mathbf{Y} = {?})$ is given by the familiar formula from elementary (i.e., non-measure-theoretic) probability theory.

If $\mathcal{G}$ is of the special form $\mathcal{G} = \sigma(I_U)$, where $I_U$ is the indicator function of an event $U\in \mathcal{G}$, then we shall write $E(X \mid U)$ in place of $E(X \mid \sigma(I_U))$.

Conditional expectations may be used to define another one of the fundamental objects in probability theory:

\begin{defn}
Let $(\Omega, \mathcal{F}, P)$ be a finite probability space and $\mathcal{G}$ a sub-$\sigma$-algebra of $\mathcal{F}$. The \textit{conditional probability of an event $U\in \mathcal{F}$ given $\mathcal{G}$}, denoted $P(U \mid \mathcal{G})$, is the $\mathcal{G}$-measurable random variable given by
	\[
	P(U \mid \mathcal{G}) = E(I_U \mid \mathcal{G}).
	\]
\end{defn}

For each atom $B\in \alpha'(\mathcal{G})$, the conditional probability $P(U\mid \mathcal{G})$ is constant on $B$ with value
	\[
	P(U\mid \mathcal{G})(B) = \frac{P(U\cap B)}{P(B)}.
	\]
For a such an atom, the mapping
	\[
	P({?}\mid \mathcal{G})(B): \mathcal{F} \to [0, \infty), \quad U \mapsto P(U\mid \mathcal{G})(B),
	\]
is a probability measure on the $\sigma$-algebra $\mathcal{F}$.

In the case that $\mathcal{G} = \sigma(\mathbf{Y})$ for some $d$-dimensional $\mathcal{F}$-measurable random vector $\mathbf{Y}$, we write $P(U \mid \mathbf{Y})$ in place of $P(U \mid \sigma(\mathbf{Y}))$. This conditional probability is a $\sigma(\mathbf{Y})$-measurable random variable, and so by the Functional Criterion for Measurability (Theorem \ref{thm:functional-criterion-measurable}) there is a function
	\[
	P(U \mid \mathbf{Y} = {?}): \mathbb{R}^d \to \mathbb{R}, \quad \mathbf{y} \mapsto P(U \mid \mathbf{Y} = \mathbf{y}),
	\]
such that
	\[
	P(U \mid \mathbf{Y} = \mathbf{y}) = P(U \mid \mathbf{Y})(\{\mathbf{Y} = \mathbf{y}\})
	\]
for those vectors $\mathbf{y}\in \mathbb{R}^d$ for which $P(\mathbf{Y} = \mathbf{y})>0$. Unpacking the definitions, we see it is given by the formula
	\[
	P(U \mid \mathbf{Y} = \mathbf{y}) = \frac{P\left(U \cap \{\mathbf{Y} = \mathbf{y}\} \right)}{P(\mathbf{Y} = \mathbf{y})}.
	\]
If the probability of the preimage $\{\mathbf{Y}=\mathbf{y}\}$ is $0$, then $P(U \mid \mathbf{Y} = \mathbf{y})=0$.

\begin{thm}\label{thm:cond-exp-familiar}
Let $X$ and $\mathbf{Y}$ be a random variable and a $d$-dimensional random vector, respectively, on a finite probability space $(\Omega, \mathcal{F}, P)$. Then for all $\mathbf{y}\in \mathbb{R}^d$, we have
	\[
	E(X \mid \mathbf{Y} = \mathbf{y}) = \int_\Omega X \, dP({?} \mid \mathbf{Y}=\mathbf{y}).
	\]
\end{thm}

\subsection{Stochastic processes}

\subsection{Martingales}

\subsection{Markov processes}

We recall from basic probability theory that a stochastic process $X_t$, defined on a probability space $(\Omega, \mathcal{F}, P)$ and indexed $t=0,1,\ldots,T$, is \textit{Markov} provided that
	\begin{equation}\label{eq:markov}
	P\left( X_{t+1} \in A \mid X_0, \ldots,X_t \right) = P \left( X_{t+1} \in A \mid X_t \right),
	\end{equation}
for all $t$ with $0\leq t < T$ and all Borel subsets $A\subset \mathbb{R}$.

A few comments about this definition are in order. First, in accordance with the general definition of conditional probabilities, \eqref{eq:markov} is an equation of \textit{random variables}. The one on the left is $\sigma(X_0,X_1,\ldots,X_t)$-measurable in general, while the one on the right is $\sigma(X_t)$-measurable. Equality of these two random variables means that they are both $\sigma(X_t)$-measurable, which is the stronger of the two measurability requirements. Second, provided that the sample space $\Omega$ is finite, there is no need to consider general Borel subsets $A$ in \eqref{eq:markov}; indeed, in this case it is equivalent to require that
	\begin{equation}\label{eq:markov-2}
	P\left( X_{t+1} = x_{t+1} \mid X_0, \ldots,X_t \right) = P \left( X_{t+1} = x_{t+1} \mid X_t \right),
	\end{equation}
for all $x_{t+1}\in \mathbb{R}$. Of course, this latter is still an equation of random variables. To verify it, it is equivalent to check the equations
	\[
	P\left( X_{t+1} = x_{t+1} \mid X_0=x_0, \ldots,X_t=x_t \right) = P \left( X_{t+1} = x_{t+1} \mid X_t=x_t \right)
	\]
of real numbers, one for each vector $(x_0,\ldots,x_{t+1})\in \mathbb{R}^{t+2}$.

The next theorem gives an equivalent characterization of Markov processes in terms of conditional expectations. We restrict to the case of a finite sample space so that the proof is easier, though the result is true in full generality.

\begin{thm}
Let $X_t$ be a stochastic process defined on a finite probability space $(\Omega, \mathcal{F}, P)$, indexed $t=0,1,\ldots,T$. Then $X_t$ is Markov if and only if
	\begin{equation}\label{eq:markov-cond}
	E\left(g(X_{t+1}) \mid X_0, X_1,\ldots,X_t\right) = E \left( g(X_{t+1}) \mid X_t \right)
	\end{equation}
for all $t$ with $0\leq t < T$ and all functions $g: \mathbb{R} \to \mathbb{R}$.
\begin{proof}
For simplicity, write
	\[
	\mathcal{F}_t = \sigma\left(X_0,X_1,\ldots,X_t \right)
	\]
for each $t$, and suppose first that the equations \eqref{eq:markov-cond} all hold. Given a Borel subset $A\subset \mathbb{R}$ and a $t$ with $0\leq t < T$, select $g=I_A$, the indicator function of $A$. Then we compute:
	\begin{align*}
	P\left(X_{t+1} \in A \mid \mathcal{F}_t\right) &= E\left( I_{\{X_{t+1} \in A\}} \mid \mathcal{F}_t \right) \\
	&= E \left( I_A(X_{t+1}) \mid \mathcal{F}_t \right) \\
	&= E \left( I_A(X_{t+1}) \mid X_t \right) \\
	&= E \left( I_{\{X_{t+1} \in A\}} \mid X_t \right) \\
	&= P\left(X_{t+1} \in A \mid X_t\right),
	\end{align*}
which is exactly what we wanted to show.

Conversely, suppose that the process $X_t$ is Markov, and fix a $t$ with $0 \leq t < T$. Since $\Omega$ is finite, there are real numbers $x_1,x_2,\ldots,x_n$ such that
	\[
	X_{t+1} = \sum_{i=1}^n x_i I_{\{X_{t+1} = x_i \}}.
	\]
If $g: \mathbb{R} \to \mathbb{R}$ is a function, then
	\[
	g(X_{t+1}) = \sum_{i=1}^n g(x_i) I_{\{X_{t+1} = x_i \}}.
	\]
We then have:
	\begin{align*}
	E \left( g(X_{t+1}) \mid \mathcal{F}_t \right) &= \sum_{i=1}^n g(x_i) E\left(I_{\{X_{t+1} = x_i \}} \mid \mathcal{F}_t \right) \\
	&= \sum_{i=1}^n g(x_i) P\left(X_{t+1} = x_i \mid \mathcal{F}_t \right) \\
	&= \sum_{i=1}^n g(x_i) P\left(X_{t+1} = x_i \mid X_t \right) \\
	&= \sum_{i=1}^n g(x_i) E\left(I_{\{X_{t+1} = x_i \}} \mid X_t \right) \\
	&= E\left( g(X_{t+1}) \mid X_t \right),
	\end{align*}
as desired.
\end{proof}
\end{thm}


\end{document}